\appendices
\section{Phụ lục tái lập kết quả}\label{sec:reproducibility}

Phụ lục này cung cấp đủ thông tin để tái lập lại thực nghiệm trong
khoảng 15 phút trên một workspace Databricks mới, theo đúng tinh thần
minh bạch và trung thực khoa học: các cấu hình chưa ghi nhận được liệt
kê rõ là \emph{chưa có bằng chứng} thay vì suy đoán.

\subsection{Mã nguồn và môi trường}

\begin{itemize}
\item \textbf{Repository:} \url{https://github.com/sonnntech/kma_hk2_chuyen_de_co_so_khoa_hoc}
\item \textbf{Nền tảng:} Databricks Free Edition; notebook Python compute hoặc serverless compute cho \texttt{notebooks/00--07}; Databricks SQL Warehouse chỉ dùng cho dashboard.
\item \textbf{Phụ thuộc cục bộ (kiểm thử):} \texttt{requirements.txt} — \texttt{pyspark>=3.5,<4.0}, \texttt{pytest>=8,<9} (Databricks Runtime đã có sẵn PySpark/Delta Lake).
\end{itemize}

\subsection{Cấu hình môi trường thực nghiệm}

\begin{table}[h]
\centering
\caption{Cấu hình môi trường (\texttt{config/demo\_config.py})}
\label{tab:repro-env}
\begin{tabular}{@{}p{3.2cm}p{4.5cm}@{}}
\toprule
\textbf{Thành phần} & \textbf{Giá trị / Nguồn} \\
\midrule
Catalog & \texttt{workspace} \\
Schema & \texttt{blockchain\_pipeline\_demo} \\
Default record count & 10\,000 \\
Random seed & 42 \\
Source system & \texttt{SCHOOL\_DEMO} \\
Benchmark record counts & 1\,000 / 5\,000 / 10\,000 / 50\,000 \\
Runs per size & 3 (+ warm-up, loại khỏi summary) \\
CPU/RAM máy chạy Databricks & \textbf{Evidence Not Found} — Databricks Free Edition/serverless compute không công bố cấu hình cố định. \\
Databricks Runtime version & \textbf{Evidence Not Found} — cần ghi lại từ workspace khi chạy lại. \\
\bottomrule
\end{tabular}
\end{table}

\subsection{Các bước tái lập (\textasciitilde15 phút)}

\begin{enumerate}
\item \textbf{Import repo:} Databricks Workspace $\to$ Repos $\to$ Add Repo $\to$ dán URL GitHub ở trên, chọn branch \texttt{main}.
\item \textbf{Chạy sạch (clean run), tuần tự bằng Python compute:}
\begin{verbatim}
notebooks/00_setup_environment.py
notebooks/01_generate_source_data.py
notebooks/02_bronze_ingestion.py
notebooks/03_silver_transformation.py
notebooks/04_gold_aggregation.py
notebooks/05_verify_integrity.py
\end{verbatim}
Kết quả kỳ vọng: \texttt{SOURCE/BRONZE/SILVER/GOLD = VALID}.
\item \textbf{Chạy một tamper scenario:} \texttt{notebooks/06\_run\_tamper\_scenarios.py} với widget \texttt{SCENARIO = MODIFY\_TRANSACTION\_AMOUNT}, \texttt{CONFIRM\_TAMPER = YES}. Reset bằng \texttt{SCENARIO = RESET\_BASELINE} trước khi đổi scenario khác.
\item \textbf{Chạy benchmark:} \texttt{notebooks/07\_experiment\_and\_metrics.py} với \texttt{RECORD\_COUNTS=1000,5000,10000,50000}, \texttt{RUNS\_PER\_SIZE=3}, \texttt{INCLUDE\_WARMUP=YES}.
\item \textbf{Xuất số liệu chính xác (khuyến nghị):} chạy KPI~8/9 trong \texttt{sql/dashboard\_queries.sql} trên \texttt{experiment\_metrics} (lọc \texttt{is\_warmup=false}) để thay các giá trị approximate trong Bảng~\ref{tab:overhead} bằng số liệu thô.
\item \textbf{Kiểm thử đơn vị cục bộ (tuỳ chọn):} \texttt{pip install -r requirements.txt \&\& pytest tests/}.
\end{enumerate}

\subsection{Known Reproducibility Gaps}

Theo đúng yêu cầu trung thực khoa học, các giới hạn sau được công khai
thay vì che giấu hoặc làm tròn số liệu:

\begin{table}[h]
\centering
\caption{Các khoảng trống về khả năng tái lập}
\label{tab:repro-gaps}
\begin{tabular}{@{}p{4cm}p{3.7cm}@{}}
\toprule
\textbf{Gap} & \textbf{Tác động} \\
\midrule
CPU/RAM và Databricks Runtime version chưa được ghi nhận & Giảm khả năng tái lập chính xác về hiệu năng; cần bổ sung khi chạy lại. \\
Bảng~\ref{tab:overhead} và Fig.~\ref{fig:overhead-chart} là giá trị approximate đọc từ dashboard chart, chưa export SQL trực tiếp & Cần chạy KPI~8/9 để lấy \texttt{avg\_overhead\_percent} và \texttt{median\_overhead\_percent} chính xác. \\
\texttt{RUNS\_PER\_SIZE = 3} & Chưa đủ mẫu cho box-plot/CDF đáng tin cậy; Fig.~\ref{fig:latency-chart} chỉ là minh họa min/điển hình/outlier, không phải phân phối đầy đủ. \\
first\_broken\_block trong Bảng~\ref{tab:results} ở mức giai đoạn tổng quát & Cần export JOIN \texttt{verification\_results} $\times$ \texttt{lineage\_events} để xác nhận chi tiết. \\
\bottomrule
\end{tabular}
\end{table}
